\section{The consistency calibration and the final theorem}
\label{sec:consistency}

\subsection{Precisely what is inherited}

Let $P_0$ be the actual ProveIt
\term{SourceRotation240.program}, and let
\[
 R_*=(3,0,0,0,3,0,0,0,0),
\]
with every higher logical register zero. The initial source
configuration has the empty path. The inherited theorem is
\begin{align}
 &\term{CalibratedRotation240.halts\_iff\_zfc\_refutation}:\notag\\
 &\hspace{5mm}
 \RHalt(P_0,(\epsilon,R_*))
 \longleftrightarrow
 \operatorname{BProv}(\mathrm{ZFCax\_s},[],\bot).
 \label{eq:calibration}
\end{align}
Its exact source is retained in
\file{reference/calibration/CalibratedRotation240.lean}.
The current repository's calibration and binary-endpoint records
identify the corresponding established proofs
\cite{calibration,rotation240}.

The right side asserts a finite derivation in the specified
first-order ZFC calculus. Finite proof trees and their numerical
codes translate by structural recursion, so it is the ordinary
arithmetized refutation statement for that presentation of ZFC.
Changing between customary effective finite-proof codings uses
the usual syntactic translations and preserves consistency.

The calibration contains substantive mathematics: correctness of
the proof enumerator and its compact calculus; correspondence to
the native theory consisting of a power-container axiom, full
global-image collection, and irreflexivity; the native theory's
comparison with ZF; and the ZF/ZFC consistency transport.
The source reduction does not re-prove these layers. They are
named inherited theorems, not inferred from a filename or from
an unqualified claim of equiconsistency.

The theorem's concrete left side is independently identified.
The source-interface checker parses all 203 literal lookup rows,
the actual lookup decision tree, the three inherited overrides,
the program root, the register-field encoding, and the logical
seed. The resulting ordered source list has exactly 6,455
triples and is the input to the new source certificate.
It is not enough merely to compare a displayed theorem name
with a separately supplied JSON file; this data check supplies
that connection.

\subsection{Direct proof of Theorem~\ref{thm:main}}

Let $P_1$ be the boundary source used by the new table.
The complete source certificate and independent interval
audit establish
\begin{equation}
 \RHalt(P_1,(\epsilon,R))
 \longleftrightarrow
 \RHalt(P_0,(\epsilon,R))
 \quad\text{for every register vector }R.
 \label{eq:sourcebisim}
\end{equation}
All instruction effects, both decrement outcomes, and the
literal halt occurrence are preserved under the same ordered
leaf indices.

The exact initialization reaches the clean frame for
$(\epsilon,R_*)$. Proposition~\ref{prop:runtime} and the
zero-padding identities of Section~\ref{sec:source} give
\begin{equation}
 \Halt(M_{238},\blank)
 \longleftrightarrow
 \RHalt(P_1,(\epsilon,R_*)).
 \label{eq:binarysource}
\end{equation}
To see the reverse implication for nonhalting explicitly,
suppose the source has no finite halt. Its next instruction is
always defined, since the source tree is finite and complete.
Each instruction has a finite positive binary realization
and ends at the next clean frame. Thus infinitely many source
steps yield unbounded binary checkpoint times. No instruction
segment contains an unmatched halt. The unique binary run
therefore cannot halt between those checkpoints.

Composing \eqref{eq:binarysource}, \eqref{eq:sourcebisim} at
$R_*$, and \eqref{eq:calibration}, followed by the finite
proof-coding equivalence, proves \eqref{eq:mainhalt}. Negating
both sides proves \eqref{eq:maincon}. Canonical renumbering is
a literal transition isomorphism, so it preserves the result.
\hfill$\square$

The local proof supplies an independent binary route:
Boundary239's own unbounded runtime has been rechecked while
certifying its guard, and Theorem~\ref{thm:fold} transfers its
halting behavior to the new table. Both routes meet at the
same source certificate and inherited calibration.

\subsection{Why the proof is formalizable in ZFC}

All new finite structures are explicit: transition tables,
source trees, affine vectors, finite words, and lists. The
new reasoning is finite case analysis, natural-number induction,
finite-list induction, elementary arithmetic, and equality
of finite-support tapes. These are ordinary constructions
within ZFC.

The inherited calibration's use of finite syntax, first-order
models, quotients, definable interpretations, and completeness
is likewise ordinary set-theoretic mathematics. Invoking that
mathematics does not require a general assertion that ZFC
proves the soundness of every theorem accepted by the whole
Lean or Rocq type theory.

The phrase \emph{provably in ZFC} here classifies the mathematical
proof. It does not report delivery of a serialized object-level
ZFC derivation or a newly checked proof-assistant theorem for
the final 238-state endpoint. The exact proof data and small
transport lemmas make those formalization goals concrete.

\subsection{Consistency and independence}

If ZFC is consistent, \eqref{eq:maincon} implies that the
machine in fact runs forever. By G\"odel's second
incompleteness theorem, consistent ZFC cannot prove that
nonhalting statement, since such a proof would prove its own
consistency.

Bare consistency should not be promoted to two-sided
independence. Excluding a proof of the false existential
halting statement requires a stronger correctness hypothesis,
such as arithmetical $\Sigma_1$ soundness. The principal
result is the requested biconditional, for which neither
an unconditional consistency assertion nor a two-sided
independence claim is needed.
